\chapter{从胰岛素到基因治疗}

\begin{kaichang}
药店的冷藏柜里有一种白色小盒子，里面躺着一支像记号笔一样的东西——胰岛素注射笔。全世界有几亿人靠它活着：每天把针头扎进肚皮，按一下按钮，几微升透明液体进入身体，血糖就有了着落。

现在请你猜三件事。第一，这一小笔液体里的胰岛素，是从哪里来的？一百年前，答案可能是"从两百头猪的胰脏里榨出来的"；今天，答案会让你意外。第二，如果一个人生病的原因是身体里的"图纸"上写错了一个字母，医生能不能不天天打针，而是直接把图纸改过来？第三，把图纸改过来和把图纸"传给孩子之前先改好"，是一回事吗？

先写下你的猜测。这一章结束时，你会发现这三个问题串起了现代医学最激动人心、也最需要冷静的一条路。
\end{kaichang}

\section{一支笔里的蛋白质}

先说你看得见的东西。得了 1 型糖尿病的人，胰岛里生产胰岛素的细胞被自己的免疫系统误伤了，胰岛素断供。胰岛素是干什么用的？它是血液里的一道"开门指令"：血糖升高时，它通知肌肉、肝脏、脂肪细胞打开细胞膜上的葡萄糖通道，把糖收进仓库。没有这道指令，糖堵在血液里，细胞却在挨饿——人会口渴、消瘦，血糖长期过高还会慢慢损坏血管、肾脏和神经。

2 型糖尿病更常见，问题不完全一样：身体还在生产胰岛素，但细胞对指令"听而不闻"，叫胰岛素抵抗。两种病的细节属于医学课，这里只要记住一件事：缺的不是糖，是一把分子钥匙。

这把钥匙本身就是一种物质。胰岛素是一个蛋白质，由 51 个氨基酸连成两条链，两条链之间用二硫键扣住，折叠成一个紧致的小球（蛋白质怎样由氨基酸折叠成机器，第 60 章讲过）。既然是物质，就有五个老问题等着我们：它由什么组成？由氨基酸。怎样连接？由基因里碱基的顺序决定。能不能变化、变多快？能——只要能找到一台会读这张图纸的机器。而生命里最擅长读图纸的机器，是细胞。

于是接下来的故事就不是"化学家怎样合成胰岛素"，而是"人怎样请别的细胞替我们生产胰岛素"。

\section{从猪胰脏到细菌工厂}

1922 年胰岛素刚被发现时，唯一的来源是屠宰场。把猪和牛的胰脏磨碎、酸提、纯化，得到一点点结晶。动物胰岛素和人胰岛素只相差一两个氨基酸，多数病人能用，但有人会过敏，而且产量始终追不上需求。

转折发生在第 82 章讲过的那套工具成熟之后。既然限制性内切酶能剪 DNA、连接酶能粘 DNA、质粒能把基因带进细菌，那么把人胰岛素的基因剪下来，粘进质粒，送进大肠杆菌——细菌每分裂一次，就把这张图纸复印一份，并且照着图纸生产人的胰岛素。一罐发酵液，就是一座微型制药厂。

\begin{zhengju}
这个故事的真实版本比传说更有意思。1978 年，一家刚成立两年的小公司基因泰克（Genentech）和洛杉矶希望之城医学中心的实验室合作，宣布用大肠杆菌生产出了人胰岛素。但请注意他们当时\emph{没有}直接从人体细胞里钓出胰岛素基因——那在技术上还很困难。他们的办法是化学合成：已知胰岛素两条链的氨基酸顺序，就反推出对应的 DNA 序列，用化学方法把两条链的基因一段段人工拼出来，分别插进两批大肠杆菌里生产，再在体外把两条链拼成完整分子。

为什么绕这么大圈子？因为"图纸顺序可以由蛋白质倒推"这件事，本身就是中心法则可信度的最好检验（第 64 章）：如果氨基酸顺序真的由碱基顺序决定，那么照着密码表反写出来的基因，放进细菌里就应该产出正确的蛋白质。实验成功了，链条闭环。1982 年，重组人胰岛素成为世界上第一个获批上市的基因工程药物。当然也有曲折：纯化出的产品必须证明和天然胰岛素一模一样，还要证明不夹带细菌的致热成分——药企用一层层色谱纯化和严格的检测回答了这些质疑。今天，工程酵母取代了大肠杆菌成为主力，产品纯度更高。
\end{zhengju}

生产出来的细菌或酵母被养在巨型不锈钢罐里，叫\textbf{生物反应器}。它不像化工反应釜那样靠高温高压，而是靠伺候细胞：温度恒定在 $37\,^{\circ}\mathrm{C}$ 左右（酵母则低一些），pH 用酸碱自动调节，营养液配方精确到每种微量元素，氧气靠搅拌和通气供应——因为细胞工厂的"工人"是活的，你虐待它们，它们就罢工甚至死亡，还会把产品降解掉。产量以克每升计，一罐几百升的发酵液，顶得上过去一个屠宰场一个月的胰脏。

\begin{qiaoliang}
感受一下数量级的变化。20 世纪 60 年代，从动物胰脏提取约 1 千克胰岛素结晶，大约需要处理十几吨胰脏——相当于十万头以上的猪。而现代酵母发酵罐生产胰岛素，产量可达每升发酵液数克。算一笔粗账：一个 1000 升的中等发酵罐，按每升 3 克计，一批就能收获约 3 千克胰岛素蛋白。也就是说，一座小型发酵车间几天的产量，相当于过去几十万头猪。这就是为什么基因工程药物一上市，动物胰岛素就迅速退场：不只是"更好"，还是"多几个数量级"。数量级的改变，常常会彻底改变一件事的社会面貌。
\end{qiaoliang}

\section{蛋白质药物的家族}

胰岛素打开了一扇门，门后是一个大家族。既然能让微生物照图纸生产蛋白质，那么任何有用的蛋白质——只要知道它的基因——原则上都能这样造。

生长激素，过去只能从去世者的脑垂体中提取，几十具遗体才够一个患儿用一年，还曾因此传播过致命的朊病毒病；换成重组产品后，这两个噩梦同时消失。凝血因子 VIII，血友病人过去靠捐献血浆浓缩制备，1980 年代曾因此被肝炎和艾滋病病毒污染；重组版本让血友病的治疗脱胎换骨。

家族里最精巧的成员是\textbf{单克隆抗体}。回忆第 70 章：抗体是免疫系统制造的 Y 形蛋白质，两臂能精确咬住特定分子（抗原）。自然免疫产生的是千百种抗体的混合物；而如果把一个产生特定抗体的 B 淋巴细胞和骨髓瘤细胞融合，得到的"杂交瘤"细胞既能无限增殖，又只生产这一种抗体——这就是单克隆抗体，1975 年的发明，后来拿了诺贝尔奖。如今用基因工程把抗体基因直接装进工程细胞，连小鼠这一步都可以省掉。单克隆抗体能咬住癌细胞表面的标记物、堵住炎症信号分子、中和病毒，是现代药物里增长最快的一类。

疫苗也在经历同样的"分子化"。传统疫苗要培养病毒再灭活或减毒，周期以年计；而新平台的思路是：不给人病毒，只给人病毒的"零件图纸"——一段编码病毒表面蛋白的 mRNA，裹在脂质纳米颗粒里送进细胞，让人体自己的细胞临时生产一小段病毒蛋白，供免疫系统演习（mRNA 怎样被翻译成蛋白质，正是第 64 章的那条链）。图纸用完即被降解，不进入细胞核，不改写基因组。2020 年新冠疫情中，这种平台把疫苗设计周期压缩到了几周——从拿到病毒基因序列开始算。

\begin{shiyan}
动手试一试：在家体验"生物反应器"的逻辑。

\textbf{材料}：干酵母一小包、白糖、三个干净矿泉水瓶、三个气球、温水、记号笔。

\textbf{步骤}：三个瓶子各装半瓶温水，各加一勺糖摇匀。甲瓶水温约 $35\,^{\circ}\mathrm{C}$（手感温热不烫），乙瓶用冰凉的自来水，丙瓶用较烫的水（约 $55\,^{\circ}\mathrm{C}$，务必请家长操作，防止烫伤）。每瓶加半勺干酵母摇匀，瓶口各套一个气球。放在原处，每隔十分钟观察一次，记录气球鼓起的情况，至少观察四十分钟。

\textbf{观察点}：哪个瓶子的气球鼓得最快？酵母里的酶把糖分解成二氧化碳和乙醇（第 58 章的发酵），气球越大说明细胞工厂开工越猛。你会发现温度有一档"恰到好处"：太冷，酶反应慢；太热，蛋白质变性，工人直接"烫死"——这就是生物反应器为什么要把温度、pH 伺候得分毫不差。

\textbf{安全提示}：全程只用食品级材料，唯一风险是热水，交给家长。做完的液体倒进下水道即可，不要饮用。
\end{shiyan}

\section{不修零件，修图纸}

以上所有药物有一个共同点：它们都是\textbf{补充}身体缺的蛋白质，而不是\textbf{修好}出错的地方。病人得打一辈子针。那么，能不能更进一步——直接修改出错的基因？

第 82 章讲了 CRISPR 等编辑工具怎样在 DNA 上定位、切割、改写。工具在手，接下来就是工程问题：把工具送进哪个细胞、送多少、结果如何。这里出现一条全书最重要的分界线，请你务必看清：

\begin{itemize}[itemsep=2pt]
\item \textbf{体细胞基因治疗}：只修改病人自己一部分普通细胞的基因。比如修改视网膜细胞治疗某种遗传性失明，修改肝脏细胞治疗某种代谢病。修改随这个人的一生结束而结束，\textbf{不会遗传给孩子}。逻辑上它和"换一批零件"类似，世界上已有多种此类疗法获批。
\item \textbf{生殖系编辑}：修改精子、卵子或早期胚胎的基因组。修改会进入这个未来孩子的\emph{每一个}细胞，包括他（她）未来的生殖细胞——改动将\textbf{代代相传}。
\end{itemize}

区别不只是技术范围，而是责任范围。体细胞治疗没治好，风险由接受了治疗的这一个病人承担，他本人在知情后做了选择；生殖系编辑的后果则由还不存在、无法同意的人承担，而且脱靶错误（第 82 章讲过，编辑工具可能切错地方）一旦写进生殖系，就像印进底版的错字，会在后代每一张复印品上重现。2018 年，有研究者宣布编辑了两个婴儿胚胎中的一个基因并使其出生，试图让她们抵抗艾滋病病毒——而那个基因的完整功能、编辑带来的脱靶损伤、以及孩子的长期健康状况全都无法保证。全球科学界几乎一致谴责了这一行为，当事人最终被判刑。目前主流共识是：生殖系编辑不具备临床应用的条件。这条界线不是反科学，恰恰是科学诚实：不知道自己工具的全部后果时，就不该在无法撤回的地方动手。

\section{把免疫细胞改造成猎人}

有一种基因治疗连病毒载体都不太用，干脆把细胞请出体外，改好再送回去——这就是 CAR-T 细胞治疗。

回忆第 70 章：T 细胞是免疫系统的猎人，靠表面的受体识别"可疑分子"，发现癌细胞就扑上去。但癌细胞很狡猾，会把自己伪装得平平无奇。CAR-T 的思路是：从病人血液里取出他自己的 T 细胞，在实验室用基因工程给它们加装一个新受体——"嵌合抗原受体"（CAR），这个受体是专门照着这种癌细胞表面的标记设计的。改造好的 T 细胞在培养皿里扩增到几亿个，再输回病人体内。等于把警察带出去特训、配上人脸识别眼镜，再放回城里。

对某些白血病和淋巴瘤，这个策略创造了奇迹：一些用尽所有疗法的病人达到了长期缓解，首批接受治疗的患儿中有人至今已无病生存超过十年。但奇迹两个字后面必须跟着账单，这一节我们一项项算：

\begin{itemize}[itemsep=2pt]
\item \textbf{疗效是有范围的}。CAR-T 目前最亮眼的战绩集中在几种血液肿瘤；对肺癌、肝癌这类实体瘤，癌细胞标记不专一、肿瘤内部环境抑制免疫细胞，效果仍然有限。
\item \textbf{风险是真实的}。几亿个被激活的免疫细胞可能同时释放大量信号分子，引发"细胞因子风暴"——高烧、低血压，严重时可危及生命；还可能误伤神经。接受 CAR-T 的病人需要在有经验的医院监护，这不是普通打针。
\item \textbf{价格与可及性}。为每个病人单独定制一批细胞，流程动辄数周，费用曾高达数百万元人民币。它目前在多数国家只用于其他疗法失败的病人，医保覆盖各地不同。
\item \textbf{时间还没有走完}。第一个 CAR-T 产品 2017 年才获批。改造过的细胞在病人体内会存续多久？几十年后有没有迟发问题？这些问题只有长期随访能回答——所以接受基因治疗的人通常会进入长达十几年的登记随访。新药获批不是故事的结局，而是观察的起点。
\end{itemize}

\begin{wuqu}
"基因治疗就是打一针，把全身的坏基因都改好，从此一劳永逸。"——两处理解都错了。第一，递送是基因治疗最大的瓶颈：你没法把编辑工具均匀地送进全身几十万亿个细胞，现有疗法大多只够得着特定器官（眼、肝、血液）里的一部分细胞，改好的细胞和没改好的细胞混在一起，叫嵌合现象。第二，"一劳永逸"忽略了生物学的不确定性：编辑可能脱靶，插入的基因可能被身体沉默，疗效可能随时间衰减。所以现实里的基因治疗更接近"一次高风险、高成本的重大维修，附带终身保修观察"，而不是魔法。
\end{wuqu}

\section{精准医学：更准，但不是算命}

把这一章串起来，你会看到一个共同信念：如果我们能读到每个人的基因（第 81 章）、能改写基因（第 82 章）、能按图纸生产蛋白质（本章），医学就能从"对所有人用同一种药"变成"按你的分子特征定制方案"。这个愿景叫\textbf{精准医学}。

它确实兑现了一部分。某些肺癌病人的肿瘤带有特定基因突变，针对该突变的药物对他们效果远好于传统化疗——先测基因再选药，已经是常规做法。同一种药，不同人代谢速度不同（肝脏里的酶有基因差异），剂量本该不同，这叫药物基因组学，第 85 章会专门讲检测能说什么、不能说什么。

但请你同时记住两条减速带。

第一条：\textbf{大多数常见病不是单基因病}。高血压、糖尿病、抑郁症背后是成百上千个基因位点各贡献一点点，再叠加饮食、运动、睡眠、运气。这类病的风险预测本质上是概率分布，不是判决。一个人带着"高风险基因型"可能终身不发病，另一个人基因平平照样中招——这和第 76 章讲的群体遗传率是同一个逻辑：统计描述的是一群人的差异，不是给某个人写好的剧本。

第二条：\textbf{证据是分等级的}。单个病例的"神奇治愈"是最弱的证据——它可能是安慰剂效应、自然病程波动或幸存者偏差。往上是病例系列、对照研究、随机对照试验，塔尖是多项随机试验的综合分析。一种新疗法从实验室到指南推荐，要沿这座塔爬很多年。看到"突破性""根治"这类词时，先问一句：证据爬到塔的哪一层了？对任何具体的治疗决定，这条标准比任何书都重要——本书只帮你理解原理，不能替代医生和正规诊疗。

\begin{tiaozhan}
用一点概率直觉看看"为什么罕见病的精准治疗也可能大量无效"。假设某种编辑疗法能成功进入目标器官 10\% 的细胞，而这 10\% 里又有 90\% 被正确修复（脱靶等失败占 10\%）。那么全身目标细胞里被修好的比例是 $0.1\times0.9=9\%$。如果恢复 20\% 的正常细胞才足以改善症状，这次治疗就达不到门槛——尽管每一步单独听起来成功率都不低。再把脱靶放进生殖系场景：假设每次编辑在某位点发生意外的概率只有万分之一，但一个胚胎的基因组有几十亿个位点、编辑工具要在多处扫描比对，意外命中一次的后验概率就不再是可以忽略的小数，而这次意外会传给所有后代。低概率事件乘以大基数，是本章所有谨慎态度的数学根源。这一段跳过不影响主线，但值得你回来重读。
\end{tiaozhan}

\section{这些疗法什么时候会失灵}

每章的老规矩：模型再好，也要说清楚它在哪里会翻车。本章这套"读写生命"的模型，有几条明确的边界。

第一，它擅长对付\textbf{单基因、已知机制、有明确目标器官}的病：一种蛋白质缺了或错了，补它或修它。但多数折磨人类的病——冠心病、阿尔茨海默病、多数癌症——是多基因、多环节、随时间演化的系统问题，没有一张单独的"错图纸"可修。把单基因病的成功外推到所有疾病，是这套模型最常见的越界。

第二，免疫系统既是帮手也是对手。身体不认识病毒载体和外来蛋白质，可能把它们当敌人攻击：轻则让治疗失效，重则引发危险的炎症。1999 年，一位名叫杰西·盖尔辛格的少年在一项早期基因治疗试验中因强烈的免疫反应去世，整个领域因此停摆整顿了好几年。这个代价换来的教训是：递送工具本身的毒性，和治疗效果一样需要当作头等科学问题，而不是工程细节。

第三，"在生产线上有效"不等于"在人体里有效"。蛋白质药怕热、怕胃酸，所以胰岛素至今只能注射不能口服——它一进消化道就被当成普通食物拆解掉了（第 63 章的消化酶可不管它身价多贵）。一种分子在培养皿里再漂亮，也要闯过吸收、分布、代谢、排泄这四关才谈得上疗效。

第四，也是最容易被忽略的一条：即使科学上全部成立，疗法还活在社会系统里。价格、产能、医保、运输冷链、地区医疗水平，决定了纸面上的突破能抵达多少人。一种只服务于少数人的神药，和一种便宜可靠的普通药，对"人类健康"这个总目标的贡献可能完全不成比例——评价医学进步时，分子精度和可及性要放在同一杆秤上。

\section{回到那支笔}

现在可以回答开场的问题了。

笔里的胰岛素从哪里来？从一座细胞工厂：人胰岛素的基因住在酵母或细菌的质粒里，发酵罐伺候着温度、pH、营养和氧气，几微升药液的背后是一条从 DNA 到蛋白质的信息流水线——你猜对了吗？

能不能不修零件、修图纸？能，但眼下只能在体细胞上、在够得着的器官里、在严格监管下做，CAR-T 和几种基因疗法已经救了一些过去无药可救的人；而生殖系编辑——把图纸"传给孩子之前先改好"——和科学界主流划出的红线重叠，那两个问题的答案不一样。

那支笔本身也在进化：胰岛素被改造成起效更快或更持久的类似物（改几个氨基酸就能改变它在皮下的聚集速度，从而改变吸收速率——结构决定性质，第 60 章的规矩在药房里照样有效），配合连续血糖监测，病人离"人工胰腺"越来越近。但从分子到病床之间，永远隔着剂量、证据、个体差异和时间。

\section{这条链通向哪里}

这一章讲的是"读写生命"在医学里的样子。同样的工具换个对象，就是完全不同的争论：把基因写进农作物，是育种还是冒险？复制一个基因组，能复制一个人吗？从零设计细胞，边界在哪里？第 84 章转入转基因、克隆与合成生物学；第 85 章则换一个角度——不改基因，只读基因，看看一张基因检测报告到底能告诉你什么、不能告诉你什么。读写生命的工具越锋利，"该不该用"这个问题就越不属于技术，而属于我们每一个人。

\section*{练一练}

\begin{sikaoti}
\item 画一张信息流图：从"人胰岛素基因"到"药店里的一支胰岛素"，标出每个环节里信息以什么形式存在（DNA？mRNA？蛋白质？），以及每个环节的"机器"是谁。图旁注明：这是表示法，罐子里并没有真的小人在读图纸。
\item 第 82 章说酶不决定反应方向、只降低活化能。把这句话用到本章：生物反应器里的细胞提供了原料、能量还是图纸？如果三者缺一，工厂会怎样？
\item 为什么从猪胰脏提取的胰岛素可能让部分病人过敏，而重组人胰岛素极少引起这个问题？请从蛋白质结构和免疫识别的角度回答（提示：第 70 章）。
\item 用你自己的话向一位长辈解释：为什么"修改我血液细胞的基因治病"和"修改胚胎的基因"是性质不同的两件事？试着不用任何专业术语。
\item 某种新药新闻里写着"治愈率高达 90\%"。列出你至少三个追问（例如：和什么相比？样本多少人？随访多久？），并说明每个问题在排除哪一种替代解释。
\item 假设一种体细胞基因治疗能把编辑工具送进肝脏 15\% 的细胞，其中 80\% 被正确修复；医生估计修好 10\% 的肝细胞即可明显改善症状。按 tiaozhan 里的方法估算，这个方案有没有希望越过门槛？如果你的结果恰好在门槛附近，你会向研究者建议先改进哪一个数字？
\end{sikaoti}
